\section{SF-E850: translating the finite ideal strengthens the non-CI obstruction}\label{SF-E850:section}

Author: /root/sixfolds, GPT-6 Astra Pro; inherited effort not exposed.
Date: 2026-11-22. Auxiliary question Q5aw: proved after author, root and independent whole-proof readings.
This is a new derivation from the accepted E849 fixed-ideal and supported-cone
identities. The four additional directions are actual translated extensions;
there is no assertion that an arbitrary extension class lifts along an
arbitrary deformation of its ideal. Author and peer statuses are recorded
separately in the accompanying manifest.

All geometric conclusions concern the auxiliary RM abelian-fourfold object
problem. The unproved transfer to the requested imaginary-quadratic Weil
eightfold remains outside this result.

\subsection{1. Precise statements}\label{SF-E850:s1}

Let \(R=\mathbf C[[x_1,x_2,x_3,x_4]]\), and let \(I\) be a proper finite-colength ideal.
Put \(A=R/I\) and

\[
 \ell=\dim_{\mathbf C}A,\qquad
 a=\dim_{\mathbf C}\operatorname{Hom}_R(I,A),\qquad
 \delta=a-3\ell,
 \qquad V=\operatorname{Ext}^1_R(I,A).
 \tag{TR.1}
\]

For \(q\ge 1\) and arbitrary \(b\) in \(\operatorname{Mat}_q(V)\), let \(M_b\) be an actual extension

\[
 0\longrightarrow A^q\longrightarrow M_b\longrightarrow I^q
       \longrightarrow0.
 \tag{TR.2}
\]

\textbf{Local theorem SF-E850a.} Without any complete-intersection, Gorenstein,
grading, scalar-matrix, or nonzero-\(b\) assumption,

\[
       \operatorname{ext}^2_R(M_b,M_b)
                    \ge 2q^2\delta+8.
 \tag{TR.3}
\]

This is a lower bound, not a claim that \(\delta\) is always nonnegative.
For the global consequence we impose the explicit sufficient hypothesis
\(\delta\ge0\).

Let \(X\) be a complex abelian fourfold and \(E\) a coherent sheaf with canonical
locally free hull and finite kernel and cokernel,

\[
 0\longrightarrow T\longrightarrow E\longrightarrow
 H=E^{\vee\vee}\longrightarrow Q\longrightarrow0,
 \qquad
 \widehat E_p\simeq M_b\oplus R^s,\quad
 b\in\operatorname{Mat}_m(V),\quad m>0,\quad s\ge0.
 \tag{TR.4}
\]

Here the actual torsion and hull quotient at \(p\) are both free of rank \(m\)
over the same \(A\); no basis is assumed equivariant. Fix a specified effective
coherent finite translation/tensor action \(G\) in supplied S21, with an actual
coherent equivariant structure on \(E\).

\textbf{Global conditional theorem SF-E850b.} Suppose \(\delta\ge0\). If a class
\(\xi\) in \(HH^2(X)\) has nonzero holomorphic bivector component and ordinary
evaluation \(\rho_E(\xi)=0\), then

\[
                 \dim\ker\sigma^2_{E,G}\ge2\delta+2>0.
 \tag{TR.5}
\]

\textbf{Balanced consequence SF-E850c.} With the genuine RM splitting and
polarization \(\theta=\theta_++\theta_-\), if

\[
 e^B\operatorname{ch}(E)
     =r(1-\nu\theta_+^2/2)(1-\nu\theta_-^2/2),
 \qquad r>0,\quad \nu>0,
 \tag{TR.6}
\]

for an algebraic rational divisor \(B\), \(E\) cannot be weakly \(G\)-semiregular.
For the identity action, (TR.4) and \(\delta\ge0\) already imply failure of
semiregularity, without (TR.6) or the evaluation hypothesis.

The strict numerical witness below has

\[
  \operatorname{Hilb}_A=(1,4,3),\qquad
  \ell=8,\qquad a=25,\qquad\delta=1,
  \qquad\operatorname{ext}^2_R(M_b,M_b)\ge2q^2+8.
 \tag{TR.7}
\]

It is not a complete intersection and does not satisfy the previous
E849 sufficient condition \(\delta\ge4\). For \(m=1\) its new invariant kernel
bound is at least four.

\subsection{2. The accepted fixed-ideal identities and their exact scope}\label{SF-E850:s2}

The elementary finite-grade calculation of E849 Section 2 gives
\(\operatorname{Hom}_R(I,R)=\operatorname{Hom}_R(I,I)=R\), \(M_b^{**}=R^q\), and

\[
  \dim V=2a-2\ell.
 \tag{TR.8}
\]

The last identity uses finite-supported duality for \(\operatorname{RHom}_R(A,A)\) and its
zero Euler characteristic; it does not assume that \(V\) is a free \(A\)-module.
For \(M=M_b\) write

\[
 \begin{aligned}
 a_L&=\dim_{\mathbf C}\{t\in\operatorname{Mat}_q(A):tb=0\},\\
 \mathcal B&=\operatorname{im}
       (\operatorname{End}_R(M)\longrightarrow\operatorname{Mat}_q(R)),
       &c&=\operatorname{length}(\operatorname{Mat}_q(R)/\mathcal B),\\
 k&=\dim\operatorname{Hom}_R(M,A^q)=q^2a+a_L,
       &\psi(t,r)&=tb-br.
 \end{aligned}
 \tag{TR.9}
\]

The map to matrices is the actual canonical hull map. Yoneda pushforward
and pullback show that the image modulo \(I\) consists exactly of those \(r\)
for which \(tb=br\) is soluble. Its dimension is \(q^2\ell-c\); projection of
\(\ker\psi\) onto the \(r\) coordinate has that image and kernel dimension \(a_L\).
Consequently E849 Section 3 proves

\[
                  \dim\operatorname{im}\psi=q^2\ell+c-a_L.
 \tag{TR.10}
\]

For clarity, the finite-hull-cone identity used below is the equal-length
case of E849 Section 4, (NI.15)--(NI.17). For \(C=\operatorname{cone}(M\to R^q)\),
all cohomology of \(\operatorname{RHom}_R(M,C)\) has finite length. If \(e_i=\operatorname{ext}^i_R(M,M)\),
its dimensions in degrees -1 through 4 are

\[
  (k, e_1+c, e_2, e_1, k+q^2\ell-c, q^2\ell).
 \tag{TR.11}
\]

Here grade-four vanishings give the first three entries, supported
duality gives the next and last entries, and the image of \(\operatorname{Hom}(R^q,M)\)
in \(\operatorname{End}(M)\), with its torsion intersection, gives the penultimate entry.
The supported complex has Euler characteristic zero because the two
defect lengths agree. Thus

\[
                  e_2=2e_1-2k+2c.
 \tag{TR.12}
\]

These are identities of finite lengths. They never apply ordinary local
self-duality to the positive-rank complex \(\operatorname{RHom}_R(M,M)\), and never
subtract infinite-dimensional endomorphism spaces.

\subsection{3. Actual translated families and the complete tangent kernel}\label{SF-E850:s3}

Let \(D=\mathbf C[\varepsilon]/(\varepsilon^2)\), \(S=R\otimes_{\mathbf C}D\). For a constant vector
\(v=(v_1,\ldots,v_4)\), let \(D_v=\sum_i v_i\partial_i\). There is a continuous
\(D\)-algebra automorphism

\[
 \phi_v:S\longrightarrow S,\qquad
 \phi_v(x_i)=x_i+\varepsilon v_i,\qquad
 \phi_v(f)=f+\varepsilon D_vf,\qquad
 \phi_v^{-1}=\phi_{-v}.
 \tag{TR.13}
\]

Continuity may be taken for the maximal ideal \((x_1,...,x_4,\varepsilon)\);
the formula on power series is also continuous for the induced formal
topology. Put \(I_v=\phi_v(I_D)\) and \(A_v=S/I_v\).

First keep \(I\) fixed and take the actual extension with class
\(b+\varepsilon\beta\), \(\beta\) in \(\operatorname{Mat}_q(V)\). Its existence and \(D\)-flatness follow
by tensoring a finite free \(R\)-resolution, exactly as in E849 (NI.13).
Apply the ring automorphism \(\phi_v\) to this entire extension. The result
is an actual \(D\)-flat \(S\)-module \(M(v,\beta)\) fitting into

\[
 0\longrightarrow A_v^q\longrightarrow M(v,\beta)
               \longrightarrow I_v^q\longrightarrow0.
 \tag{TR.14}
\]

Both the ideal and the extension class have been transported. This
construction does not require a map from Ext at the old ideal onto
Ext at an independently chosen deformed ideal.

The actual relative canonical hull is \(S^q\), its image is \(I_v^q\), and its
finite quotient is \(A_v^q\), compatibly with the special fiber. Indeed
tensoring the free resolution first gives \(\operatorname{Ext}^j_S(A_D,S)=0\) for \(j<4\);
applying \(\phi_v\) preserves those vanishings. Applying \(\operatorname{Hom}(-,S)\) to
(TR.14) and to \(0\to I_v\to S\to A_v\to 0\) then proves all of these identifications
through the canonical evaluation maps. This is a direct calculation
over the nonreduced \(S\), not an assumed double-dual base-change theorem.

The family defines a linear tangent map

\[
 \mathcal L:\mathbf C^4\oplus\operatorname{Mat}_q(V)
                  \longrightarrow\operatorname{Ext}^1_R(M,M).
 \tag{TR.15}
\]

Linearity follows either by the first-order extension description of
module deformations, or by choosing finite presentations and cocycle
representatives: every coefficient changes by the sum of its derivative
in \(v\) and its \(\beta\) term; mixed terms contain \(\varepsilon^2\) and vanish.
No higher-order or noncommutative deformation is invoked.

We now determine the whole kernel. If \(\mathcal L(v,\beta)=0\), there is an \(S\)-linear
isomorphism \(M(v,\beta)\simeq M\otimes_{\mathbf C}D\) inducing the identity on the special
fiber. This is precisely the zero class in \(\operatorname{Ext}^1_R(M,M)\). Functoriality
of the canonical evaluation map gives an \(S\)-linear isomorphism of their
relative hull quotients. Because \(q>0\),

\[
       I_v=\operatorname{Ann}_S(A_v^q)
          =\operatorname{Ann}_S(A_D^q)=I_D.
 \tag{TR.16}
\]

This uses the actual quotient under an arbitrary middle-module
isomorphism. It does not suppose the isomorphism preserves a chosen
presentation or fixes a chosen basis of the hull.

Equation (TR.16) implies \(D_v(I)\subset I\). If \(v\) is nonzero, choose a linear
form \(l\) with \(D_v(l)=1\). The derivation descends to the finite-dimensional
algebra \(A\). Write \(m_l\) for multiplication by \(l\) in \(A\). The derivation rule
gives

\[
             [D_v,m_l]=I_A,\qquad
             0=\operatorname{tr}[D_v,m_l]=\ell,
 \tag{TR.17}
\]

which is impossible in characteristic zero because \(\ell>0\). Hence \(v=0\).

With \(v=0\), the exact fixed-ideal kernel calculation E849 (NI.12)--(NI.14)
applies. A trivialization induces canonical kernel and hull
automorphisms \(1+\varepsilon t\) and \(1+\varepsilon r\). Their extension equality is
\((1+\varepsilon t)(b+\varepsilon\beta)=b(1+\varepsilon\bar r)\), so \(\beta=b\bar r-tb\).
Conversely every such pair gives an actual extension isomorphism, and
every matrix over \(A\) lifts to \(R\). Therefore

\[
 \ker\mathcal L=\{0\}\oplus\operatorname{im}\psi,\qquad
 e_1\ge4+q^2(2a-3\ell)-c+a_L.
 \tag{TR.18}
\]

Substituting (TR.18) and \(k=q^2a+a_L\) into the supported identity (TR.12)
proves (TR.3). The extra eight comes from four actual independent
translation directions, counted twice by that identity.

\subsection{4. Rectangular stabilizer blocks, including zero ranks}\label{SF-E850:s4}

We give the actual finite-action interface so that the added directions
are not incorrectly required to be invariant deformations.

At a point stabilizer of a translation/tensor action, the translation
is identity. Flat local frames for the torsion Picard factors give
one constant projective scalar cocycle. The completed module has the
actual twisted-Maschke decomposition

\[
                 \widehat E_p=\bigoplus_\lambda
                               S_\lambda\otimes N_\lambda.
 \tag{TR.19}
\]

The stabilizer acts on \(S_{\lambda}\) and is \(R\)-linear; it does not act by
nontrivial formal coordinate transformations. Apply the exact
multiplicity functor to the canonical torsion, hull, image, and
quotient. Its idempotents are \(A\)-linear on the finite modules and
\(R\)-linear on the hull. Direct summands of \(A^m\) and \(R^{m+s}\) are free,
since both rings are local. Write the torsion and quotient ranks
\(t_{\lambda}\) and \(q_{\lambda}\), and hull rank \(h_{\lambda}\).

A surjection \(R^{h_{\lambda}}\to A^{q_{\lambda}}\) has a unit minor modulo the
maximal ideal. Ordinary lifted row and column operations put it into
the residue map on \(q_{\lambda}\) coordinates and zero on the rest.
The kernel is \(I^{q_{\lambda}}\oplus R^{h_{\lambda}-q_{\lambda}}\). The extension over
the free summand splits. Thus

\[
 N_\lambda\simeq M_\lambda\oplus R^{h_\lambda-q_\lambda},\qquad
 0\longrightarrow A^{t_\lambda}\longrightarrow M_\lambda
                 \longrightarrow I^{q_\lambda}\longrightarrow0.
 \tag{TR.20}
\]

The assumed ordinary evaluation vanishing restricts to every actual
summand. On the completed affine neighborhood, the higher HKR
components vanish; in abelian logarithm coordinates only the nonzero
constant bivector \(\pi\) remains. Consequently \(\rho_{\pi}(N_{\lambda})=0\).

The accepted E819 local finite-cone argument, with the exact scope
rechecked in E849 Section 5, applies to each \(N_{\lambda}\). Its first-order
free-resolution lift and hull-map lift produce a finite perfect cone
over the dual numbers for the algebra with product
\(a*b=ab+\varepsilon\pi(da,db)/2\). The coordinate commutator has zero
supertrace on this cone. For a nonzero \(\pi_{ij}\) this gives

\[
            0=\varepsilon\pi_{ij}(q_\lambda-t_\lambda)\ell,
            \qquad t_\lambda=q_\lambda.
 \tag{TR.21}
\]

No balanced-character hypothesis on an individual \(N_{\lambda}\) was needed.
If \(q_{\lambda}=0\), this also forces \(t_{\lambda}=0\). Some \(q_{\lambda}>0\) because
the original quotient \(A^m\) is nonzero. For such a block the new square
theorem applies. Its translation families are ordinary deformations
used to bound ordinary Ext dimensions, not asserted equivariant
families. The classes \(\operatorname{Id}_{S_\lambda}\otimes u\) in its \(\operatorname{Ext}^2\) are nevertheless
actual stabilizer invariants, because the stabilizer acts on \(S_{\lambda}\)
only. This is the same faithful invariant subspace used by E849.

\subsection{5. The actual global injection and the six-dimensional trace bound}\label{SF-E850:s5}

Let \(C_X=\operatorname{cone}(E\to H)\). Its cohomology has finite support. All cohomology
sheaves of \(\operatorname{RHom}(E,C_X)\) also have finite support, so restriction and
completion identify every global \(\operatorname{Ext}^j(E,C_X)\) with the direct sum of
the corresponding completed local groups.

Grade-four finite-defect vanishings give local \(\operatorname{Ext}^1(E,H)=\operatorname{Ext}^2(E,H)=0\).
The preceding global map \(\operatorname{Ext}^1_X(E,H)\to \operatorname{Ext}^1_X(E,C_X)\) is zero: the
target restriction is the finite-support isomorphism, whereas the
source restriction followed by the local map is zero. Therefore the
global hull triangle gives the canonical injection

\[
 \bigoplus_{p'}\operatorname{Ext}^2_{\widehat{\mathcal O}_{X,p'}}
                  (\widehat E_{p'},\widehat E_{p'})
                  \hookrightarrow\operatorname{Ext}^2_X(E,E).
 \tag{TR.22}
\]

This is the actual map of E849 (NI.22)--(NI.23), not an inference from
localized vanishing alone. Its Serre pairing against evaluation of
\(\gamma\) in \(HH^2\) is the trace of a connecting product in \(\operatorname{Ext}^3(E,C_X)\).
Finite support makes that product depend only on completed restriction
of the evaluation. Thus only \(H^0(X,\bigwedge^2T_X)\), of dimension six,
contributes. Full CY4 trace/evaluation adjunction bounds semiregularity
on every subspace of (TR.22) by rank six.

Insert \(\operatorname{Id}_{S_\lambda}\otimes\operatorname{Ext}^2(M_\lambda,M_\lambda)\) into the block at
\(p\) and transport it around its actual orbit. Stabilizer invariance
makes the sum well-defined; its component at \(p\) proves injectivity.
The canonical map (TR.22) is equivariant. Supplied S21 identifies the
ordinary invariant map with the relevant identity-sector map. Hence

\[
 \dim\ker\sigma^2_{E,G}
    \ge 2q_\lambda^2\delta+8-6
    \ge2\delta+2>0,\qquad \delta\ge0.
 \tag{TR.23}
\]

This proves E850b. For the identity action one uses the given square
block directly, so the evaluation and rectangular reduction are unnecessary.

For E850c, the entire balanced character has the two cap-kernel
classes \(\pi_\pm+\nu\eta_\pm\), with nonzero bivectors and
\(\eta_{\pm}=\pi_{\pm}\) contracted with \(\theta_\pm^2/2\). Algebraic B-field
conjugation preserves their bivector parts. Published BF composition
and supplied S21 equivariant compatibility imply that hypothetical
weak injectivity forces the ordinary evaluation of either class to
vanish. E850b then contradicts that injectivity.

The exact source/dependency locators are unchanged from E849 Sections
2,5,6: Stacks 47.18 (tag0A81), Theorem 47.18.3/Lemma 47.18.4;\cite{SFStacksE847LocalDuality} Lipman
Lemma 7.3/Theorem 7.4, pp.~64--69;\cite{SFLipman1984} Stacks 15.76.2 (tag 0658);\cite{SF-StacksPerfect} Stacks48.27,
Lemma 48.27.1(5),(7), Remarks 48.27.2--3 and Lemma 48.27.4;\cite{SFStacksE846Duality} published
BF2008 Proposition 6.2.1/Theorem 6.4.2;\cite{SF-BF2008} and supplied DEG\_L S20's first
proof paragraph and S21.6 through the adjunction before S21.8.
No new whole upstream source audit is claimed here. The new proof is
the actual translation family, annihilator kernel, and strict witness.

\subsection{6. An explicit length-eight algebra with \(\delta=1\)}\label{SF-E850:s6}

Write \(x,y,z,w\) for the four variables. Let \(I\) be generated by

\[
 \begin{gathered}
 xz,\quad xw,\quad yw,\quad xy-yz,\quad yz-zw,\\
 x^2-2y^2+z^2,\qquad y^2-2z^2+w^2.
 \end{gathered}
 \tag{TR.24}
\]

Here is both an exact algebra model and a finite rank certificate.
Let \(A\) have vector-space basis \(1,x,y,z,w,s_1,s_2,s_3\). Its maximal
ideal cubed is zero, the \(s_k\) annihilate the maximal ideal, and the
coefficients of \(s_1,s_2,s_3\) in \(x_ix_j\) are the entries of

\[
 B_1=I_4,\qquad B_2=\operatorname{diag}(1,2,3,4),\qquad
 B_3=\begin{pmatrix}
 0&1&0&0\\1&0&1&0\\0&1&0&1\\0&0&1&0
 \end{pmatrix}.
 \tag{TR.25}
\]

This multiplication is associative: every triple of maximal-ideal
elements is zero. The seven quadrics (TR.24) are linearly independent
and vanish in this model. Multiplying them by the four variables gives
a 28-by-20 cubic coefficient matrix. The accompanying exact certificate
exhibits a nonzero 20-by-20 minor modulo 101, so these products span all
cubics over \(\mathbf C\). Hence \(I\) contains \(\mathfrak m^3\) and the model is exactly \(R/I\), with
Hilbert function (1,4,3). The three \(s_k\) form its socle: a linear
combination of \(x_i\) annihilated by every \(x_j\) has all coefficients zero
already by \(B_1=I_4\). There are seven independent minimal generators,
so it cannot be a complete intersection in a regular local ring of
dimension four. This is invariant under formal coordinates and changes
of ideal generators.

There are at least 25 independent homomorphisms \(I\to A\). Every linear map
\(I_2\to A_2\) extends to an \(R\)-linear map by killing \(I_{\ge3}\); multiplication
by the maximal ideal kills \(A_2\). These give \(7\cdot3=21\) degree-zero maps.
The four partial derivatives \(I\to A\) give four maps of degree minus one.
Their independence is the constant-derivation trace argument (TR.17),
or the directly checked rank-four derivative matrix. They are independent
of the first 21 maps because their quadratic-generator values have
different degrees. Thus \(a\ge25\).

For the reverse bound, let \(M_i\) be the 8-by-8 multiplication matrices
of \(x_i\) in the displayed basis. The Koszul complex on the four commuting
endomorphisms \(x_i-M_i\) of \(R\otimes_{\mathbf C}A\) is a free \(R\)-resolution of \(A\).
One direct verification translates a power series by the commuting
nilpotent matrices \(M_i\); this invertible operation conjugates the four
operators to the ordinary coordinate operators. All joint products
of three \(M_i\) vanish, so the matrix translation is a finite Taylor
expression. Applying \(\operatorname{Hom}_R(-,A)\) yields the commutator-Koszul complex
with first maps

\[
 \begin{aligned}
 d_0:\operatorname{End}_{\mathbf C}(A)&\longrightarrow
                      \operatorname{End}_{\mathbf C}(A)^4,
             &F&\longmapsto([M_i,F])_i,\\
 d_1:\operatorname{End}_{\mathbf C}(A)^4&\longrightarrow
                      \operatorname{End}_{\mathbf C}(A)^6,
             &(F_i)&\longmapsto([M_i,F_j]-[M_j,F_i])_{i<j}.
 \end{aligned}
 \tag{TR.26}
\]

Their integer matrix shapes are 256-by-64 and 384-by-256. \(\operatorname{End}_R(A)=A\),
because \(A\) is cyclic, so \(\operatorname{rank}_{\mathbf C}(d_0)=64-8=56\). The exact certificate
gives a 175-by-175 submatrix of \(d_1\) with nonzero determinant modulo 101.
It follows that \(\operatorname{rank}_{\mathbf C}(d_1)\ge175\) and

\[
 a=\operatorname{ext}^1_R(A,A)
       =256-\operatorname{rank}(d_1)-56\le25.
 \tag{TR.27}
\]

The lower bound proves equality. In particular \(\delta=25-3\cdot8=1\).
The certificate separately evaluates the chosen minors; it does not
treat a modular rank upper bound as a characteristic-zero lower bound
on Ext. The lower bound came from the 25 explicit maps. The result
is (TR.7).

The complete checker is \texttt{\detokenize{rm_nonci_translation_certificate.py}}; its
output includes the 175 selected rows/columns and determinant residue,
the cubic minor, all 512 associativity checks, all six multiplication
commutators, \(d_1d_0=0\), the socle rank, and the derivative rank.
The complete first execution passed; its actual timestamps and exact
scope are retained separately. The preliminary four-algebra exploration
is retained and is not described as a universal search or proof.

\subsection{7. Actual coherent objects and the exact separation claim}\label{SF-E850:s7}

Take the accepted rank 32 balanced bundle \(H_0\) on the explicit CM RM
fourfold of E815/E820. At a point \(p\) embed \(\operatorname{Spec}A\) with the local
equations (TR.24), and choose a surjection \(H_0\to\mathcal O_Z\). Let \(J\) be its
kernel. Any prescribed \(b\) in \(\operatorname{Ext}^1_R(I,A)\) globalizes to an actual
extension \(0\to\mathcal O_Z\to E_b\to J\to0\): finite support gives the global/local
Ext identification, while \(\operatorname{Ext}^j_X(H_0,\mathcal O_Z)=0\) for \(j>0\). Locally
the actual module is \(M_b\oplus R^{31}\). Its canonical hull is \(H_0\) and
the two finite defects are both \(\mathcal O_Z\). Therefore

\[
                   [E_b]=[H_0],\qquad
                   \operatorname{ch}(E_b)=\operatorname{ch}(H_0).
 \tag{TR.28}
\]

There are nonzero choices because \(\dim V=2a-2\ell=34\). For \(b\) nonzero,
the rank-one module \(M_b\) is indecomposable: a decomposition would have
a rank-zero summand, necessarily a finite direct summand of its torsion
\(A\). The local cyclic module \(A\) is indecomposable, so that summand would
be all of \(A\), splitting the extension and forcing \(b=0\).

These are actual coherent and hence ordinary gluable objects. They
are not semiregular seeds. No nontrivial finite equivariant structure
is manufactured; the theorem rules out every specified effective
coherent action that actually exists, and the identity action always
exists.

The strict comparison is precise. The common local algebra is not CI,
so the earlier CI-specific Artin and matrix criteria do not apply to
it. Its \(\delta=1\) fails E849's stated \(\delta\ge4\) hypothesis. In the
rank-one common-defect case, E849's bound alone supplies only two
local middle-Ext dimensions against a possible rank-six trace; the
new actual translation directions supply ten and hence a nonzero
kernel. This expands the stated sufficient local criterion, without
asserting that every previous theorem in the project has been classified
or that the auxiliary-to-target-eightfold bridge has been proved.

The remaining scope includes algebras with \(\delta<0\), unequal/nonfree
defect modules over a common algebra, different algebras, singular hulls,
and more general derived objects. No assertion of existence in those
remaining cases is made.
